% Preâmbulo comum às duas versões do artigo no formato SBC Reviews: main.tex (português) e
% main-en.tex (inglês). Cada versão define antes os textos que dependem do idioma:
% \SBCrotulopendente, \SBCorcidpendente, \SBCemailpendente, \versaoconsultada e \publicacaoantecipada.
% A classe sbcreviews-2025.cls não é alterada; ver TEMPLATE-SBC-REVIEWS.md.
\usepackage[bottom]{footmisc}
\usepackage{afterpage}
\usepackage{url}
% O estilo apalike-sol.bst imprime o DOI como texto; o pacote permite o sublinhado de um DOI fora do modo matemático.
\usepackage[strings]{underscore}
\setcitestyle{square}

\definecolor{engtitle}{rgb}{0.5,0.5,0.5}
\definecolor{darkblue}{rgb}{0.0,0.0,0.0}
\hypersetup{colorlinks,breaklinks,
            linkcolor=darkblue,urlcolor=darkblue,
            anchorcolor=darkblue,citecolor=darkblue}

% ---------- Pendências visíveis ----------
\newcommand{\pendencia}[1]{\par\noindent\textcolor{red!65!black}{\textbf{\SBCrotulopendente:} #1}\par}
\newcommand{\aconfirmar}[1]{\textcolor{red!65!black}{#1}}

% ---------- Referências ----------
% Identificador e endereço das referências que só existem no arXiv (campo howpublished do .bib).
\newcommand{\soartigo}[1]{#1}

% ---------- Adaptação dos arquivos escritos para a classe IEEEtran ----------
% Os capítulos abrem os apêndices com \appendices, comando da IEEEtran.
\newcommand{\appendices}{\appendix}
% As tabelas foram dimensionadas para a classe IEEEtran (coluna de 252 pt e página de 516 pt); aqui a
% coluna tem 240 pt e a página, 498 pt. Sem alterar os arquivos dos capítulos, o espaço entre colunas
% diminui e a tabela que ainda passar da largura disponível é reduzida até caber (no máximo 5\%).
\newsavebox{\SBCcaixa}
\newcommand{\SBCtabelasajustadas}{%
  \let\SBCtabular\tabular \let\SBCendtabular\endtabular
  \renewenvironment{tabular}%
    {\begin{lrbox}{\SBCcaixa}\SBCtabular}%
    {\SBCendtabular\end{lrbox}%
     \ifdim\wd\SBCcaixa>\linewidth\resizebox{\linewidth}{!}{\usebox{\SBCcaixa}}\else\usebox{\SBCcaixa}\fi}}
\AtBeginEnvironment{table}{\setlength{\tabcolsep}{4.5pt}\SBCtabelasajustadas}
\AtBeginEnvironment{table*}{\SBCtabelasajustadas}

% ---------- Campos do cabeçalho da revista: marcadores do template, preenchidos pelos editores ----------
\jid{ROCS}
\jtitle{SBC Computing Reviews, 202X, XX:1}
\issn{2966-3938}
\doi{10.5753/rocs.202X.XXXXXX}
\copyrightstatement{This work is licensed under a Creative Commons Attribution 4.0 International License}
\jyear{202X}

% ---------- Título e autoria ----------
\newcommand{\titulotcc}{Revisão Rápida da Literatura: Engenharia de Prompt para Codificação e Implementação de Software}
% Título em inglês proposto por Claude Code em 24/11/2026; pendente de aprovação do autor (D141).
\newcommand{\titulotccen}{Prompt Engineering for Software Coding and Implementation: A Rapid Literature Review}
\newcommand{\nomeorientador}{Eduardo Magno Lages Figueiredo}
\newcommand{\nomecoorientador}{Pedro Saint Clair Garcia}

% Autoria: orientador como coautor por decisão do autor em 08/11/2026 (D88); coorientador informado
% em 21/11/2026 (D136), coautor por pedido do autor em 22/11/2026 (D138), com vínculo ao DCC/UFMG
% informado pelo autor (D139), não conferido em fonte institucional.
% ORCID e e-mail vêm de busca em fontes públicas feita em 24/11/2026 (D142) e estão A CONFIRMAR:
%  - Eduardo Figueiredo: ORCID no registro público do ORCID (vínculo UFMG, doutorado em Lancaster);
%    e-mail no perfil do DCC/UFMG (https://dcc.ufmg.br/professor/eduardo-figueiredo/).
%  - Pedro Garcia: ORCID de «Pedro Garcia» (UFMG), com obra assinada por «Pedro S. C. Garcia» em
%    coautoria com Eduardo Figueiredo (doi:10.5753/sbes.2024.3503). E-mail não localizado.
%  - Caio César de Freitas: ORCID não localizado; e-mail institucional não informado.
% O envelope marca o autor correspondente; a escolha do primeiro autor também está a confirmar.
\author[Freitas et al. 202X]{
\affil{\textbf{Caio César de Freitas}~\aconfirmar{[\SBCorcidpendente]}~\textcolor{blue}{\faEnvelope}~~[{Universidade Federal de Minas Gerais}~|~\aconfirmar{\textit{\SBCemailpendente}}~]}

\affil{\textbf{\nomeorientador}~\orcidlink{0000-0002-6004-2718}~[{Universidade Federal de Minas Gerais}~|\href{mailto:figueiredo@dcc.ufmg.br}{~{\textit{figueiredo@dcc.ufmg.br}}}~]}

\affil{\textbf{\nomecoorientador}~\orcidlink{0009-0005-7744-697X}~[{Universidade Federal de Minas Gerais}~|~\aconfirmar{\textit{\SBCemailpendente}}~]}
}
